\section{Problem Formulation}
\label{sec:problem}

\paragraph{Quote--pay--observe setting.}
We consider a single autonomous buyer with initial wallet budget $B$ that repeatedly chooses among $K$ payable service providers over rounds $t \in \{0,\ldots,T-1\}$. At the start of round $t$, the buyer observes a task context $\mathbf{x}_t \in \mathcal{X}$ and the wallet state $B_t$. A programmable payment interface exposes enough price information to determine which providers are affordable before payment; let $p_{a,t}$ denote the pre-payment quote for provider $a$, and define
\begin{equation}
\label{eq:affordable-set}
\mathcal{A}_t = \{a \in \{1,\ldots,K\}: p_{a,t} \leq B_t\}.
\end{equation}
The quote is available before selection; the realized charged cost $c_t$ is observed after the selected payment and is kept as the cost feedback signal across rounds. If $\mathcal{A}_t$ is empty, the run stops. Otherwise, the buyer selects one provider $a_t \in \mathcal{A}_t$, commits to the paid call through the protocol, and cannot observe the service outcome before paying.

\paragraph{Feedback and wallet dynamics.}
After the selected call completes, the buyer observes only the chosen provider's outcome:
\[
(q_t, c_t, f_t) \sim \mathcal{D}_{a_t,t}(\cdot \mid \mathbf{x}_t),
\]
where $q_t \in [0,1]$ is a task-level quality score, $c_t \in \mathbb{R}_{\geq 0}$ is the realized charged cost, and $f_t \in \{0,1\}$ marks timeout, parse failure, or another non-delivery event. Outcomes and realized charges for unchosen providers are not observed as samples; this is the bandit feedback that makes provider choice a learning problem rather than a full-information ranking problem. The wallet evolves deterministically as
\begin{equation}
\label{eq:wallet}
B_{t+1} = B_t - c_t, \qquad B_0 = B,
\end{equation}
and $\tau \leq T$ denotes the first round at which the horizon is reached or no provider is affordable.

\paragraph{Utility and objective.}
We separate intrinsic service delivery from payment pressure. Utility combines quality with a structural failure penalty:
\begin{equation}
\label{eq:utility}
\ut = q_t - \nu f_t.
\end{equation}
The penalty captures retry overhead, broken call chains, and downstream agent stalls that are not fully represented by setting $q_t=0$ on failure.

The payment-aware reward combines utility with normalized realized cost:
\begin{equation}
\label{eq:reward}
\ctilde = \mathrm{clip}(c_t/c_{\max},0,1), \qquad
\lamn = \frac{\lambda_t}{1+\lambda_t}, \qquad
\rt = (1-\lamn)\ut - \lamn\ctilde .
\end{equation}
Here $c_{\max}$ is a fixed cost normalizer and $\lambda_t \geq 0$ is the wallet-pressure multiplier defined below. The bounded weight $\lamn \in [0,1]$ keeps utility and cost comparable across wallet states while preserving the sign of failures. A policy $\pi$ aims to maximize the expected cumulative payment-aware reward
\begin{equation}
\label{eq:objective}
J(\pi) = \mathbb{E}_{\pi}\!\left[\sum_{t=0}^{\tau-1} \rt \right],
\end{equation}
subject to the affordability constraint in Eq.~\eqref{eq:affordable-set}. In PA-DCT, beliefs are updated from the two observed channels $(\ut,c_t)$ rather than from the scalar $\rt$: $\ut$ tracks provider quality, $c_t$ tracks realized payment cost, and $\rt$ defines the decision/evaluation objective.

\paragraph{Budget pressure.}
The wallet-pressure multiplier is a deterministic function of cumulative spend. For $t>0$, define the realized burn rate as the fraction of the wallet spent per elapsed round,
\begin{equation}
\label{eq:burn-rate}
\mathrm{actual\_burn\_rate}_t = \frac{(B-B_t)/B}{t},
\qquad
\mathrm{target\_rate} = \frac{1}{T}.
\end{equation}
The normalized burn deviation is
\begin{equation}
\label{eq:lambda}
\bex_t =
\frac{\mathrm{actual\_burn\_rate}_t-\mathrm{target\_rate}}
{\mathrm{target\_rate}},
\qquad
\lambda_t =
\begin{cases}
1, & t=0,\\
\exp(\alpha \bex_t), & t>0,
\end{cases}
\end{equation}
with sharpness parameter $\alpha \geq 0$. Thus $\lamn=0.5$ when the buyer is exactly on the linear burn plan, decreases when the buyer is under-spending, and increases toward cost-dominated decisions when the wallet is being spent too quickly.

\paragraph{Non-stationarity and scope.}
Provider distributions $\mathcal{D}_{a,t}(\cdot \mid \mathbf{x})$ may change over time in utility, reliability, or realized cost. Change points are not announced to the policy. Tasks arrive from an exogenous workload stream, and providers are non-strategic posted-price endpoints: they expose quotes and deliver services rather than bidding strategically for allocation. Reverse auctions, multi-agent coordination, dispute resolution, and latency-aware routing are outside this formulation.

\paragraph{Why a contextual bandit.}
The only controlled action in a round is the provider choice $a_t$; the immediate feedback is the selected call's $(\ut,c_t)$, and the only persistent state affected by the action is deterministic wallet accounting. There is no delayed reward, within-round planning horizon, or latent transition model to learn. We therefore formulate agent-native payment decision-making as a contextual bandit with budget state and non-stationary feedback, rather than as a full reinforcement-learning problem.
